\begin{problem}[Closed points versus rational points]
Compare the closed points and base-field rational points of
\[
\Spec\CC[t],\qquad
\Spec\RR[t],\qquad
\Spec\FF_p[t].
\]
Explain algebraically why the three pictures differ.
\end{problem}

\paragraph{Why this problem matters.}
It prevents one of the most persistent confusions in algebraic geometry: closed points are
not always coordinate tuples over the chosen base field.

\paragraph{Useful ideas before starting.}
Closed points of $k[t]$ correspond to irreducible polynomials.  A point is $k$-rational
when the quotient by its maximal ideal is $k$.

\paragraph{Guided hints.}
Over $\CC$ all irreducibles are linear; over $\RR$ they are linear or irreducible
quadratic; over $\FF_p$ they occur in every degree.

\begin{solution}
For $\CC[t]$, every irreducible polynomial is $t-a$, so every closed point is
\[
(t-a)
\]
and has quotient $\CC$.  Thus every closed point is $\CC$-rational.

For $\RR[t]$, closed points are generated by irreducible real polynomials.  The linear
ones $(t-a)$ have quotient $\RR$ and are rational.  An irreducible quadratic $q$ gives
\[
\RR[t]/(q)\cong\CC,
\]
so it is a closed point but not an $\RR$-rational point.

For $\FF_p[t]$, each irreducible polynomial $f$ of degree $d$ gives a closed point with
quotient
\[
\FF_p[t]/(f)\cong\FF_{p^d}.
\]
Only $d=1$ gives quotient $\FF_p$, so only degree-one closed points are
$\FF_p$-rational.
\end{solution}

\paragraph{What happened mathematically.}
The topological spectrum naturally remembers algebraic extensions of the base field.

\paragraph{Extensions.}
Relate this to the weak Nullstellensatz for finitely generated algebras over an
algebraically closed field.

\paragraph{Provenance.}
New synthesis problem collecting field-sensitive examples scattered through legacy
Algebraic Geometry 1 and Commutative Algebra 13.
